% =====================================================================
%  Gemeinsame Präambel für alle Zusammenfassungen und Spickzettel
%  Einbinden mit:  \input{../../_latex/praeambel}
% =====================================================================

% --- Sprache und Schrift ---------------------------------------------
\usepackage[ngerman]{babel}
\usepackage[T1]{fontenc}
\usepackage{lmodern}
\usepackage{microtype}
% In Schreibmaschinenschrift keine Ligaturen: "--" bleibt zwei Striche, "<<" bleibt "<<"
\DisableLigatures{encoding = *, family = tt*}
\usepackage{csquotes}

% --- Seitenlayout (Ränder im Hauptdokument mit \geometry{} anpassen) --
\usepackage[a4paper]{geometry}

% --- Mathematik ------------------------------------------------------
\usepackage{amsmath, amssymb, mathtools}

% --- Tabellen und Listen ---------------------------------------------
\usepackage{array, booktabs, tabularx}
\usepackage{enumitem}

% --- Farben, Boxen, Grafiken -----------------------------------------
\usepackage[table]{xcolor}
\usepackage[most]{tcolorbox}
\usepackage{graphicx}
\graphicspath{{bilder/}}
\usepackage{tikz}
\usetikzlibrary{arrows.meta, positioning, calc}
\usepackage{forest}            % Bäume zeichnen (BST, Heap, AVL)

% --- Programmcode ----------------------------------------------------
\usepackage{listings}

% --- Links (immer zuletzt laden) -------------------------------------
\usepackage[hidelinks]{hyperref}

% =====================================================================
%  Farben
% =====================================================================
\definecolor{hsluBlau}{HTML}{1F4E79}
\definecolor{defFarbe}{HTML}{1F4E79}
\definecolor{merkFarbe}{HTML}{C55A11}
\definecolor{bspFarbe}{HTML}{548235}
\definecolor{pruefFarbe}{HTML}{A50021}
\definecolor{codeHintergrund}{HTML}{F5F5F5}

% =====================================================================
%  Boxen
%    \begin{definition}{Titel} ... \end{definition}
%    \begin{merke} ... \end{merke}
%    \begin{beispiel}{Titel} ... \end{beispiel}
%    \begin{pruefung} ... \end{pruefung}      (Prüfungstipp)
% =====================================================================
\tcbset{
  basis/.style={enhanced, breakable, boxrule=0.6pt, arc=2pt,
                left=6pt, right=6pt, top=4pt, bottom=4pt,
                fonttitle=\bfseries, before skip=8pt, after skip=8pt}
}
\newtcolorbox{definition}[1]{basis, colframe=defFarbe, colback=defFarbe!4,
  title={Definition: #1}}
\newtcolorbox{merke}{basis, colframe=merkFarbe, colback=merkFarbe!5,
  title={Merke}}
\newtcolorbox{beispiel}[1]{basis, colframe=bspFarbe, colback=bspFarbe!4,
  title={Beispiel: #1}}
\newtcolorbox{pruefung}{basis, colframe=pruefFarbe, colback=pruefFarbe!4,
  title={Prüfungstipp}}

% =====================================================================
%  Wochenziele-Tabelle
%    \begin{wochenziele}
%      \ziel{Lernpunkt}{Lernziel}
%    \end{wochenziele}
%  Die Nummern werden automatisch gezählt.
% =====================================================================
\newcounter{zielnr}
\newcommand{\ziel}[2]{\stepcounter{zielnr}\thezielnr & #1 & #2 \\}
\newenvironment{wochenziele}{%
  \setcounter{zielnr}{0}%
  \par\noindent
  \renewcommand{\arraystretch}{1.35}%
  \rowcolors{2}{hsluBlau!5}{white}%
  \tabularx{\linewidth}{@{}r>{\raggedright\arraybackslash}p{0.3\linewidth}>{\raggedright\arraybackslash}X@{}}
  \toprule
  \textbf{Nr.} & \textbf{Lernpunkt} & \textbf{Lernziel} \\
  \midrule
}{%
  \bottomrule
  \endtabularx
  \par
}

% =====================================================================
%  Programmcode
%    \begin{lstlisting}[language=Python] ... \end{lstlisting}
% =====================================================================
\lstset{
  basicstyle=\ttfamily\small,
  backgroundcolor=\color{codeHintergrund},
  keywordstyle=\color{hsluBlau}\bfseries,
  commentstyle=\color{gray}\itshape,
  stringstyle=\color{bspFarbe},
  numbers=left, numberstyle=\tiny\color{gray},
  frame=single, rulecolor=\color{gray!40},
  breaklines=true, showstringspaces=false, tabsize=4,
  literate={ä}{{\"a}}1 {ö}{{\"o}}1 {ü}{{\"u}}1
           {Ä}{{\"A}}1 {Ö}{{\"O}}1 {Ü}{{\"U}}1 {ß}{{\ss}}1
}

% =====================================================================
%  Kurzbefehle
% =====================================================================
\newcommand{\Oh}[1]{\mathcal{O}\!\left(#1\right)}   % \Oh{n^2}
\newcommand{\N}{\mathbb{N}}
\newcommand{\Z}{\mathbb{Z}}
\newcommand{\R}{\mathbb{R}}
\renewcommand{\vec}[1]{\mathbf{#1}}                 % Vektoren fett
